\documentclass{article}

% NeurIPS 2026 style - submission mode (anonymous, with line numbers)
\usepackage[position]{neurips_2026}

\usepackage[utf8]{inputenc}
\usepackage{amsmath,amssymb,amsthm}
\usepackage{booktabs}
\usepackage{graphicx}
\usepackage{hyperref}
\usepackage{xcolor}
\usepackage{multirow}

\newtheorem{theorem}{Theorem}
\newtheorem{definition}{Definition}
\newtheorem{finding}{Finding}

\title{Knowledge Graph Uncertainty Research Should Adopt Scope Conditions}

\author{}

\begin{document}

\maketitle

\begin{abstract}
\textbf{This position paper argues that knowledge graph (KG) uncertainty research lacks scope conditions---explicit statements of when methods are valid.} We present evidence that the relationship between coverage and reliability \emph{inverts} depending on KG structure: coverage predicts reliability in temporal/hierarchical KGs (AUROC $\approx$ 0.70) but \emph{anti-predicts} it in multi-relational KGs (AUROC $\approx$ 0.40). Across 22 benchmarks, we observe a striking 13-vs-0 structural split. We argue this is not merely a methodological oversight but reflects deeper assumptions the field has inherited from supervised learning---assumptions that do not transfer to the relational setting. We call for uncertainty research to adopt scope conditions: explicit characterizations of when methods apply. This paper is intended to open discussion, not close it; we highlight open questions that the community must address.
\end{abstract}

\section{Introduction}

Uncertainty quantification (UQ) for knowledge graph embeddings is an active research area with compelling motivations. KGE models are deployed in drug discovery \citep{bonner2022review}, fraud detection, and recommendation systems where knowing \emph{when not to trust} a prediction is critical.

Yet something is wrong with how the field evaluates uncertainty methods.

\textbf{The core problem:} Methods are developed on one benchmark, evaluated on another, and deployed on a third---with no theory of when results should transfer. The implicit assumption is that uncertainty behaves similarly across knowledge graphs. We argue this assumption is false, and worse, that the field lacks the conceptual vocabulary to even ask whether it holds.

\begin{center}
\fbox{\parbox{0.92\textwidth}{
\textbf{Central Position:} KG uncertainty research needs \emph{scope conditions}---explicit statements of the structural properties under which a method provides valid uncertainty estimates. Without scope conditions, we cannot interpret cross-benchmark results, aggregate metrics, or deployment decisions. The coverage-reliability inversion we document is a symptom of this deeper problem.
}}
\end{center}

\subsection{The Evidence in Brief}

We analyzed 22 diverse KG benchmarks and found a striking pattern: the relationship between coverage (whether entities have appeared with a relation in training) and prediction reliability \emph{inverts} depending on structural type.

\begin{itemize}
    \item \textbf{13/13 multi-relational KGs} show partial coverage \emph{outperforming} full coverage
    \item \textbf{0/9 structured KGs} (temporal, hierarchical) show this pattern
\end{itemize}

This is not a subtle effect. On FB15k-237, partial coverage achieves 59.5\% Hits@10 versus 32.3\% for full---nearly double. A practitioner using coverage to triage predictions could, under certain conditions, flag the \emph{more} reliable ones.

But we do not present this finding as the main contribution. The inversion is a \emph{symptom}. The underlying issue is the absence of scope conditions in uncertainty research.

\section{How Did We Get Here? The Inherited Assumptions}

Why does KG uncertainty research lack scope conditions? We argue it inherited assumptions from supervised learning that do not transfer to the relational setting.

\subsection{The IID Assumption and Its Relational Failure}

In supervised learning, uncertainty methods assume training and test data are drawn from the same distribution. This assumption, while often violated, provides a clear scope condition: the method is valid when IID holds. KG uncertainty research inherits a related assumption implicitly: prediction reliability correlates with training exposure. This assumption is rarely stated explicitly but underlies common practices---flagging predictions involving unseen entity-relation pairs as uncertain, or trusting predictions for frequently-observed patterns.

Knowledge graphs break this assumption fundamentally. A test triple $(h, r, t)$ may involve:
\begin{itemize}
    \item Entities seen in training with the \emph{same} relation (full coverage)
    \item Entities seen with \emph{different} relations (partial coverage)
    \item Entities never seen with the relation (zero coverage)
\end{itemize}

These are structurally different prediction problems, yet uncertainty methods treat them identically. The field inherited the IID framing without asking whether it applies.

\subsection{The Implicit Coverage Assumption}

A deeper assumption underlies coverage-based uncertainty: \emph{more training signal implies more reliable predictions}. This seems obviously true---until you examine it.

In multi-relational KGs, high-coverage entities participate in \emph{many different relations}. Their embeddings must satisfy conflicting constraints. A gene that is the subject of ``encodes'' but object of ``associated with disease'' must represent both roles in a single vector. Full coverage may indicate not reliable knowledge, but \emph{representational overload}.

In temporal KGs, entities have consistent roles across relations. An actor who ``visits'' also ``negotiates with'' and ``cooperates with''---similar contexts that reinforce rather than conflict. Here, coverage does indicate reliable knowledge.

The coverage assumption holds in one structural regime and inverts in another. But because the assumption was never made explicit, we never asked when it holds.

\subsection{The Benchmark Monoculture}

The field's benchmark choices reinforced these blind spots. FB15k-237 and WN18RR became the default evaluation pair. But these datasets have opposite structural properties:
\begin{itemize}
    \item \textbf{FB15k-237}: Multi-relational, high context divergence, coverage misleads
    \item \textbf{WN18RR}: Hierarchical, low context divergence, coverage helps
\end{itemize}

A method achieving ``good results on both'' might simply be averaging over a structural discontinuity. Aggregate metrics hide rather than reveal the problem.

\section{The Evidence: A Structural Split}

\begin{table}[t]
\centering
\caption{Coverage-reliability relationship across 22 KG benchmarks. Multi-relational KGs (encyclopedic + biomedical) show partial coverage outperforming full coverage in all 13 cases (AUROC $<$ 0.5). Structured KGs (temporal, hierarchical, dense) show the expected pattern in all 9 cases (AUROC $>$ 0.6). We use ComplEx \citep{trouillon2016complex}; whether this extends to other architectures is an open question.}
\label{tab:main}
\small
\begin{tabular}{llcccc}
\toprule
\textbf{Type} & \textbf{Dataset} & \textbf{Full} & \textbf{Partial} & \textbf{Paradox?} & \textbf{AUROC} \\
\midrule
Encyclopedic & FB15k-237 & 32.3\% & \textbf{59.5\%} & \textcolor{red}{Yes} & 0.40 \\
& CoDEx-M & 32.0\% & \textbf{59.0\%} & \textcolor{red}{Yes} & 0.39 \\
& YAGO3-10 & 12.0\% & \textbf{44.2\%} & \textcolor{red}{Yes} & 0.39 \\
& ConceptNet & 0.0\% & \textbf{0.7\%} & \textcolor{red}{Yes} & 0.41 \\
& NELL & 2.8\% & \textbf{3.6\%} & \textcolor{red}{Yes} & 0.40 \\
\midrule
Biomedical & DRKG & 0.6\% & \textbf{1.7\%} & \textcolor{red}{Yes} & 0.40 \\
& PrimeKG & 1.9\% & \textbf{2.9\%} & \textcolor{red}{Yes} & 0.40 \\
& Hetionet & 1.5\% & \textbf{2.3\%} & \textcolor{red}{Yes} & 0.39 \\
\midrule
Temporal & ICEWS14 & \textbf{59.0\%} & 18.0\% & No & 0.72 \\
& ICEWS18 & \textbf{42.0\%} & 15.0\% & No & 0.74 \\
& GDELT & \textbf{35.0\%} & 10.0\% & No & 0.71 \\
\midrule
Hierarchical & WN18RR & \textbf{45.0\%} & 7.0\% & No & 0.71 \\
& WN18 & \textbf{48.0\%} & 12.0\% & No & 0.68 \\
\midrule
Dense & Kinship & \textbf{10.4\%} & 5.0\% & No & 0.65 \\
\bottomrule
\end{tabular}
\end{table}

\noindent{\footnotesize Full results for all 22 datasets in supplementary material.}

Table~\ref{tab:main} presents representative results. The pattern is consistent: multi-relational KGs show the paradox (13/13), structured KGs do not (0/9). Fisher's exact test: $p < 10^{-4}$.

We propose \textbf{context divergence} as a candidate structural discriminant:
\begin{equation}
D(e) = 1 - \frac{|R_h(e) \cap R_t(e)|}{|R_h(e) \cup R_t(e)|}
\end{equation}
where $R_h(e)$ and $R_t(e)$ are relations where entity $e$ appears as head and tail. Multi-relational KGs have mean divergence $\approx$ 0.77; structured KGs have $\approx$ 0.52. Whether this is the \emph{causal} factor or merely correlated remains an open question.

\section{What Should Change: Scope Conditions for Uncertainty}

\subsection{The Concept of Scope Conditions}

We borrow the concept of \emph{scope conditions} from social science methodology \citep{walker1967scope}. A scope condition specifies the boundary conditions under which a theoretical claim holds. ``X causes Y'' becomes ``X causes Y \emph{when Z holds}.''

Uncertainty research currently lacks scope conditions. Claims like ``energy scores provide reliable uncertainty'' or ``coverage indicates prediction quality'' are stated without specifying when they apply.

\subsection{Proposed Reforms}

\textbf{Immediate:} Report metrics separately for multi-relational and structured KGs. Make structural type explicit in evaluation.

\textbf{Medium-term:} Develop diagnostic tools (e.g., context divergence) that practitioners can compute before deployment.

\textbf{Long-term:} Build a theory of structural validity---for which structural properties does method X provide valid uncertainty estimates?

\section{Open Questions for the Community}

We deliberately leave key questions unresolved:

\textbf{1. Is context divergence causal or merely correlated?}

Our data show divergence correlates with the paradox, but we have not proven causation. An intervention study---artificially modifying divergence and observing effects---would strengthen the claim. We have not done this.

\textbf{2. Does the split extend beyond ComplEx?}

All our experiments use ComplEx, a bilinear model. TransE, DistMult, and RotatE share similar embedding structure and likely exhibit the same pattern. But GNN-based models (R-GCN, CompGCN) have different inductive biases. Does the split hold? We do not know.

\textbf{3. Can training objectives eliminate the paradox?}

Perhaps modified loss functions or regularization could prevent the embedding interference that causes the paradox in multi-relational KGs. No existing method achieves this, but that does not mean it is impossible.

\textbf{4. What other scope conditions exist?}

We identified structural type as one confounding variable. There may be others: relation arity, graph density, entity type distributions. A systematic study of scope conditions is needed.

\textbf{5. Should coverage-based UQ be abandoned entirely?}

Our data suggest coverage should not be trusted for multi-relational KGs. But should the field abandon it entirely, or adapt it with structure-aware methods?

\section{Alternative Views}

\textbf{``This is just an empirical finding, not a position.''}

Fair criticism. We present evidence, then interpret it through the lens of scope conditions. Whether this framing is useful is for the community to judge.

\textbf{``22 datasets is not enough.''}

True. The 13-vs-0 split is striking but not definitive. We may have missed edge cases or structural subtypes. More data is needed.

\textbf{``Practitioners already evaluate on their target domain.''}

Hopefully. But published methods are often presented as general-purpose, and deployment decisions may rely on benchmark results from structurally incompatible datasets. Our concern is with how the field communicates validity.

\section{Related Work}

\textbf{KG uncertainty.} UKGE \citep{chen2019embedding}, BEUrRE \citep{chen2021probabilistic}, and calibration approaches \citep{guo2017calibration} address KGE uncertainty without scope conditions. We hypothesize they exhibit structure-dependent behavior.

\textbf{KG structural analysis.} \citet{akrami2020realistic} and \citet{rossi2021knowledge} characterize structural properties. We extend this by arguing structure determines uncertainty validity.

\textbf{Scope conditions in ML.} The ML community increasingly recognizes that methods have implicit scope conditions \citep{lake2018generalization}. We apply this lens to KG uncertainty specifically.

\section{Conclusion}

We have argued that KG uncertainty research needs scope conditions: explicit statements of when methods are valid. The coverage-reliability inversion we document---present in all 13 multi-relational KGs, absent in all 9 structured KGs---is evidence that current methods have implicit, unstated scope conditions.

This paper is intended to open discussion. We have presented evidence, proposed a framework (scope conditions), and highlighted open questions. We do not claim to have solved the problem; we claim the problem exists and deserves attention.

\textbf{The core message:} Before claiming an uncertainty method ``works,'' we should ask: works \emph{for which KGs}? Until we can answer this, uncertainty research findings are difficult to interpret.

\bibliographystyle{plainnat}
\bibliography{references}

\end{document}
